\documentclass[11pt]{article}
\usepackage[a4paper,margin=28mm]{geometry}
\usepackage{amsmath,amssymb,amsthm,booktabs,hyperref}
\newcommand{\Q}{\mathbb Q}
\newcommand{\Z}{\mathbb Z}
\newcommand{\rad}{\operatorname{rad}}
\newtheorem{theorem}{Theorem}
\newtheorem{lemma}[theorem]{Lemma}
\newtheorem{proposition}[theorem]{Proposition}
\newtheorem{corollary}[theorem]{Corollary}
\title{Parity-sensitive height bounds for mixed powers}
\author{}
\date{16 November 2026}
\begin{document}
\maketitle

\begin{abstract}
For rational $x,y$ and an integer $n\ge3$, a fixed weak rational
$abc$ inequality bounds the heights of $x$ and $y$ in terms of the
height of $x^{2n}+y^2$. The constants are explicit in the exponent
and constant of that inequality. The estimate for $x^{16}+y^2$
also gives a height budget for a two-argument detector whose residual
divisor has degree four in its rational input. Under ordinary $abc$
and the eight-contact estimate, the corresponding absorption coefficient
is $1406/2065<1$.
\end{abstract}

For a reduced rational number $a/b$, with $b>0$, put
$h(a/b)=\log\max(|a|,b)$ and $h(0)=0$. For a positive integer $N$,
$\rad(N)$ is the product of its distinct prime divisors.

\section{A fixed-exponent height inequality}

Assume the following weak rational $abc$ inequality for fixed real
numbers $\kappa\ge1$ and $K\ge1$:
\begin{equation}\label{eq:weakabc}
 C\le K\rad(ABC)^\kappa
 \quad\text{if }A+B=C,\quad
 A,B,C\in\Z_{>0}\text{ are pairwise coprime}.
\end{equation}

\begin{theorem}\label{thm:mixed}
Let $n\ge3$ be an integer and suppose that
$1\le\kappa<2n/(n+2)$. Set
\[
 D=2n-\kappa(n+2)>0,\qquad
 c_{n,\kappa}=\frac{n+\kappa(n-2)}{nD}.
\]
Under~\eqref{eq:weakabc}, for every $x,y\in\Q$ and
$s=x^{2n}+y^2$,
\begin{align}
 h(x)&\le c_{n,\kappa}h(s)+\frac{\log K}{D},\label{eq:xheight}\\
 h(y)&\le n c_{n,\kappa}h(s)+\frac{n\log K}{D}.
 \label{eq:yheight}
\end{align}
\end{theorem}

\begin{proof}
If $xy=0$, the assertions follow from the height of a power, since
$c_{n,\kappa}\ge1/(2n)$ and $n c_{n,\kappa}\ge1/2$.
This includes $x=y=0$.

Otherwise write $s=M/N>0$ in lowest terms. Reducing the positive
ratio $x^{2n}/y^2$ gives coprime positive integers $A,B$; with
$C=A+B$ we have
\begin{equation}\label{eq:primitive}
 x^{2n}=s\frac AC,\qquad y^2=s\frac BC.
\end{equation}
The ratio $A/B=(x^n/y)^2$ is a rational square. Since $A,B$ are
coprime, both are integer squares. Also $A,B,C$ are pairwise coprime.

We claim the exact radical majorant
\begin{equation}\label{eq:majorant}
 \rad(ABC)^{2n}\le A B^n C M^{2n-1}N^{n-1}.
\end{equation}
At a prime let $a=v(A)$, $b=v(B)$, $c=v(C)$ and $e=v(s)$.
At most one of $a,b,c$ is positive, $a,b$ are even, and
\begin{equation}\label{eq:congruences}
 a-c+e\equiv0\pmod{2n},\qquad b-c+e\equiv0\pmod2.
\end{equation}
The required inequality at that prime is
\[
 2n\,\boldsymbol1_{p\mid ABC}
 \le a+nb+c+(2n-1)\max(e,0)+(n-1)\max(-e,0).
\]
Here are all cases with $p\mid ABC$.
If $a>0,e\ge0$, then $a+e$ is a positive multiple of $2n$.
If $a>0,e<0$, then $a\ge2$ and $-e\ge2$ by the second
congruence in~\eqref{eq:congruences}; hence
$a+(n-1)(-e)\ge2n$.
If $b>0$, then $b\ge2$ and $nb\ge2n$.
If $c>0,e=0$, then $c\ge2n$ by the first congruence.
If $c>0,e>0$, then $c+(2n-1)e\ge2n$.
Finally, if $c>0,e<0$, then $c-e$ is a positive multiple of
$2n$, and
$c+(n-1)(-e)\ge c-e\ge2n$.
Primes outside $ABC$ contribute zero on the left and a nonnegative
quantity on the right. This proves~\eqref{eq:majorant}.

Put $H=\log C$, $m=\log M$ and $d=\log N$. Since $A,B\le C$,
\[
 \log\rad(ABC)
 \le\frac{(n+2)H+(2n-1)m+(n-1)d}{2n}.
\]
Applying~\eqref{eq:weakabc} and rearranging gives
\begin{align*}
 H&\le
 \frac{\kappa((2n-1)m+(n-1)d)+2n\log K}{D}\\
 &\le\frac{\kappa(3n-2)h(s)+2n\log K}{D}.
\end{align*}
The product-height inequality and~\eqref{eq:primitive} yield
$2n h(x)\le h(s)+H$ and $2h(y)\le h(s)+H$.
Substitution, together with
$D+\kappa(3n-2)=2(n+\kappa(n-2))$,
proves both bounds.
\end{proof}

\begin{lemma}\label{lem:shift}
If $s\in\Q_{\ge0}$ and $B\in\Z_{\ge1}$, then
$h(s)\le h(B+s)$ and $h(B+s)\ge\log B$.
\end{lemma}
\begin{proof}
Write $s=M/N$ in lowest terms with $M\ge0,N>0$.
The fraction $(M+BN)/N$ is reduced, and its numerator is at least
$M$, at least $N$, and at least $BN$.
\end{proof}

For $\kappa=101/100$, some exact constants in
Theorem~\ref{thm:mixed} are as follows. The last column is the
coefficient of $\log K$ in~\eqref{eq:xheight}.
\begin{center}
\begin{tabular}{ccc}
\toprule
$n$ & $c_{n,\kappa}$ & $1/D$\\
\midrule
4 & $301/388$ & $50/97$\\
6 & $251/588$ & $25/98$\\
8 & $703/2360$ & $10/59$\\
\bottomrule
\end{tabular}
\end{center}

\begin{corollary}\label{cor:16}
Assume~\eqref{eq:weakabc} with $\kappa=101/100$. For
$X,T\in\Q$, $B\in\Z_{\ge1}$ and $a=B+X^{16}+T^2$,
\[
 h(X)\le\frac{703}{2360}h(a)+\frac{10}{59}\log K,
 \qquad
 h(T)\le\frac{703}{295}h(a)+\frac{80}{59}\log K,
 \qquad h(a)\ge\log B.
\]
All constants are independent of $B,X,T$.
\end{corollary}

\section{An abstract detector height budget}

\begin{proposition}\label{prop:budget}
Fix $n,\kappa$ as in Theorem~\ref{thm:mixed}, positive constants
$\gamma,d,A_0$, and fixed nonnegative error constants.
Suppose a collection of rational witnesses, depending on an integer
$B\ge1$, has quantities $X,T,t_0,t_1\in\Q$,
$m\in\Z$ and $N\in\Z_{\ge1}$ satisfying
\begin{align*}
 t_1-t_0&=a=B+X^{2n}+T^2,\\
 h(t_j)&\le\gamma\log N+C_0\qquad(j=0,1),\\
 \log N&\le d h(X)+\log(1+|m|)+C_1,\\
 m^2&\le A_0h(T)+J_0.
\end{align*}
Assume that these constants, and the constant $K$ in
\eqref{eq:weakabc}, are independent of $B$ and the witnesses.
If
\begin{equation}\label{eq:budget}
                         2\gamma d c_{n,\kappa}<1,
\end{equation}
then $B$ is bounded above by a fixed constant. Thus such witnesses
do not exist for all sufficiently large fixed $B$.
\end{proposition}
\begin{proof}
The difference-height inequality gives
$h(a)\le2\gamma\log N+2C_0+\log2$.
Theorem~\ref{thm:mixed} and Lemma~\ref{lem:shift} therefore give
\[
 h(X)\le2\gamma c_{n,\kappa}
       \bigl(dh(X)+\log(1+|m|)\bigr)+C_2.
\]
Condition~\eqref{eq:budget} absorbs the $h(X)$ term. Consequently
$h(X)\ll\log(1+|m|)+1$, then $\log N$ and $h(a)$ satisfy
the same bound. The companion bound for $h(T)$ gives
$h(T)\ll\log(1+|m|)+1$. Combining this with the assumed quadratic
bound on $m$ bounds $|m|$, and hence $h(a)$, uniformly. Finally
$\log B\le h(a)$ bounds $B$. Every implied constant is independent
of $B$ and the witnesses.
\end{proof}

\section{Eight contacts and a product residual}

For fixed distinct rational contacts, the marked-contact consequence
of ordinary rational $abc$ gives
\[
 h(t)\le\frac1{r-4-2\delta}\log M(t)+O(1)
 \quad(r>4+2\delta),
\]
where $r$ is the number of contacts and $M(t)$ is their positive
odd-valuation radical outside a fixed exceptional set; see
\cite[Section~2]{Marked}. For eight contacts and $\delta=1/4$ the
coefficient is $2/7$. A contact argument itself has bounded height
and can be included by assigning its radical the value one.
This marked-contact assertion uses ordinary $abc$; the single fixed
exponent in Theorem~\ref{thm:mixed} alone is not being used to assert it.

Let $X=u/v$ be reduced with $v>0$, let $\beta=X^2+3$, and let
$n_0=m/2\in\tfrac12\Z$. If $n_0\ne\beta$, put
\begin{equation}\label{eq:productresidual}
 N=(u^2+3v^2)\left|mv^2-2u^2-6v^2\right|>0.
\end{equation}
At every odd prime where $\beta$ is integral and
$\beta(n_0-\beta)\equiv0\pmod p$, one has $p\nmid v$ and
$p\mid N$. Also
\begin{equation}\label{eq:Nbound}
                    N\le32(1+|m|)e^{4h(X)}.
\end{equation}
Indeed the two factors in~\eqref{eq:productresidual} are bounded by
$4e^{2h(X)}$ and $8(1+|m|)e^{2h(X)}$, respectively.

\begin{corollary}\label{cor:product}
Assume ordinary rational $abc$. Suppose a two-argument detector has
the product-residual divisibility just described at every consumed
prime, uses $a=t_1-t_0=B+X^{16}+T^2$, and supplies the all-point
index bound $m^2\le A_0h(T)+J_0$ with fixed constants.
Then, for all sufficiently large fixed $B$, it cannot have witnesses
with $n_0\ne X^2+3$.
\end{corollary}
\begin{proof}
Each noncontact radical divides~\eqref{eq:productresidual}.
The contact case has bounded height. Hence
$h(t_j)\le(2/7)\log N+O(1)$ uniformly. Apply
Proposition~\ref{prop:budget} with $n=8$, $\kappa=101/100$,
$\gamma=2/7$ and $d=4$. Its strict coefficient is
\[
        \frac47\,4\,\frac{703}{2360}
                 =\frac{1406}{2065}<1.
\]
\end{proof}

If the conclusion $n_0=X^2+3$ is reached, then $2X^2=m-6\in\Z$;
the reduced denominator of $X$ has square dividing two, so $X\in\Z$.
The corollary is a height implication. Its hypotheses include the
detector's local divisibility and index theorem; it does not construct
such a detector or supply its rational completeness theorem.

The required scale has a quadratic presentation using the shared
coordinates
\[
 z_2=X^2,\qquad z_4=z_2^2,\qquad z_8=z_4^2,
 \qquad t_1-t_0=B+z_8^2+T^2.
\]
Relative to the presentation of $B+X^8+T^2$ using $z_2,z_4$, this
adds exactly one rational coordinate and one quadratic equation.
Therefore a replacement detector satisfying the hypotheses of
Corollary~\ref{cor:product} gives a strict coordinate saving precisely
when its other components save at least two rational coordinates after
all squares, height coordinates and nonzero constraints are included.
This count concerns the height conversion and makes no assertion that
such a replacement has been obtained.

\begin{thebibliography}{9}
\bibitem{Marked}
\emph{A quartic in 78 variables from a mixed-power height estimate},
project companion,
\nolinkurl{papers/quartic_existence/quartic_abc78_20261115_v34.tex},
Sections~2--3.
\end{thebibliography}
\end{document}
